%% Methods of Information in Medicine -- Original Article (Methods Print track).
%% Budget: first 5 printed pages free (~22,000 characters incl. spaces, minus 1,500
%% per figure/table, references included); EUR 125 per extra page. Agreed ceiling:
%% 10 printed pages (EUR 625), ~44,000 characters incl. floats; aim for ~8.
%% Check with ./budget.py.
%% Final submission must be .docx (pandoc), double-spaced, line numbers, figures as
%% separate TIFF/EPS files, black and white only (colour costs EUR 350).
%% The earlier SWJ draft (22 July 2026) was replaced by this one; a copy is in
%% ../../workspace/swj-draft-2026-07-22.zip.
\documentclass[12pt]{article}
\usepackage[utf8]{inputenc}
\usepackage[T1]{fontenc}
\usepackage[margin=1in]{geometry}
\usepackage{setspace}
\usepackage{lineno}
\usepackage{booktabs}
\usepackage{graphicx}
\usepackage{url}
\usepackage[super,sort&compress,comma]{natbib}
\usepackage{xcolor}
\usepackage{hyperref}

\newcommand{\code}[1]{\texttt{#1}}
%% Values still to be filled from eval/results/. Never type a number by hand.
\newcommand{\TBD}[1]{\textcolor{red}{[#1]}}

\doublespacing
\linenumbers

\title{Composing HealthDCAT-AP Dataset Descriptions from REDCap Projects and
Machine-Actionable Data Management Plans}
\author{\TBD{Authors, degrees, departments, affiliations; corresponding author's
mailing address, telephone and e-mail}}
\date{}

\begin{document}
\maketitle

\section*{Abstract}

\paragraph{Background} The European Health Data Space requires health data holders
to describe their datasets with a fixed set of metadata elements, expressed in the
HealthDCAT-AP profile, from March 2029. Much clinical research data is collected in
REDCap projects, whose data dictionaries already describe the variables. A REDCap
project does not record how its dataset is governed, and that information is usually
held in the study's data management plan.

\paragraph{Objectives} To develop a method that composes a HealthDCAT-AP dataset
description from a REDCap project, a machine-actionable data management plan and a
set of institutional settings, with terminology codes attached to the dataset, its
forms and its fields, and to determine which of the legally required metadata elements
each source can supply.

\paragraph{Methods} We implemented the method as an open-source tool that reads a REDCap
data dictionary or project export without patient records, reads a plan in the Research
Data Alliance common standard, and emits a HealthDCAT-AP record with a variable-level
description in CSV on the Web. Terminology codes are carried in REDCap field annotations.
We built descriptions for a clinical research project in three cumulative stages, ran
the tool over public REDCap projects, and validated the output against the official
SHACL shapes.

\paragraph{Results} \TBD{Qualitative summary once eval/results are in: share of
the required elements supplied by REDCap, by institutional settings and by the plan;
elements that only the REDCap programming interface or a person can supply;
conformance outcome; robustness over the public projects.}

\paragraph{Conclusions} \TBD{One or two sentences, qualitative, from the results.}

\paragraph{Keywords} Metadata; Electronic Data Capture; Data Management; Systematized
Nomenclature of Medicine; Semantic Web
%% MeSH check: "Metadata" (D000081085?) "Electronic Data Capture"? -> VERIFY each term
%% in MeSH Browser before submission; max 5.

\newpage
\section{Introduction}

Regulation (EU) 2025/327 on the European Health Data Space (EHDS) requires each health
data access body to publish a catalogue of the datasets available for secondary use
(Article~77).\citep{ehds2025} Implementing Regulation (EU) 2026/2098 sets the minimum
metadata elements that health data holders must provide for these catalogues. It
requires the elements to be expressed with the definitions and controlled vocabularies of
HealthDCAT-AP, the health extension of the DCAT application profile for data portals in
Europe, and it applies from 26 March 2029.\citep{ir20262098,healthdcatap,dcatap} Every
dataset that a research institute holds and that falls under the regulation will
therefore need a HealthDCAT-AP description.

In academic clinical research a large share of these datasets is collected with REDCap,
an electronic data capture platform used by institutions in many countries.%
\citep{harris2009redcap,harris2019redcap} Each REDCap project carries a data dictionary
that defines its forms, its fields, their types and their permitted values. The variable
description that HealthDCAT-AP requires can be derived from that dictionary. Most of the
remaining elements describe how the dataset is governed: the publisher, the person
responsible, the access conditions, the legal basis and whether the data are personal.
REDCap does not record this information. A data management plan (DMP) does, and when the
plan follows the Research Data Alliance (RDA) DMP Common Standard it can be read by
software.\citep{rdadmp,miksa2019ten}

Current tools address parts of this task. The guideline of the TEHDAS2 joint action
explains to data holders how to describe a dataset in HealthDCAT-AP,\citep{tehdas2d51}
and an online editor lets a person enter the description by hand.\citep{hipeditor}
ROME, a REDCap module, annotates fields and answer options with ontology concepts while a
project is designed,\citep{meigen2026rome} and RareLink defines REDCap data dictionaries
for rare-disease registries with ontology codes in their field annotations.%
\citep{graefe2025rarelink} The \code{sdc\_healthdcat\_ap} project derives HealthDCAT-AP
records from published data models of the Semantic Data Charter.\citep{sdchealthdcat}
To our knowledge, no published method produces a HealthDCAT-AP description from a data
capture project and completes it with the governance metadata of a machine-actionable
plan. \TBD{State how the related-work search was done: databases, terms, date.}

\section{Objectives}

The objectives of this work were:
(1) to define a method that composes a HealthDCAT-AP dataset description from three
sources, namely a REDCap project, a machine-actionable DMP and a set of institutional
settings;
(2) to attach terminology codes to the description at the level of the dataset, the
form and the field, with a convention for carrying the codes in REDCap;
(3) to determine which of the minimum metadata elements of Implementing Regulation
2026/2098 each source can supply; and
(4) to test the conformance of the output against the official shapes and the robustness
of the method over public REDCap projects.

\section{Methods}

\subsection{Sources and their roles}

The method combines three sources (Figure~\ref{fig:architecture}). The REDCap project
supplies the title, the forms, the fields with their types and value sets, and
structural facts such as longitudinal events. The institutional settings supply values
that are the same for every project of an institution, such as the publisher, the
default access rights and the health data category; they are written once in a
configuration file. The DMP supplies the values that differ between studies and are
decided outside the data capture system: the description, the creator and contact
point, the licence, the funding project and the declaration of personal data. A REDCap
project and a plan rarely carry the same value. Where they do, the value from REDCap is
kept and the plan fills only empty elements. The institutional settings fill the
elements that remain empty after both.

The method reads two REDCap file formats: the data dictionary in comma-separated values
(CSV) and the project export in XML, an extension of the CDISC Operational Data Model
(ODM).\citep{cdiscodm} It reads only the project's structure. An export that contains
patient records is rejected with an error, and a separate command removes the records
from such an export so that its structure can be described. The tool never writes to
REDCap and runs on the user's computer, so no project file is sent to a remote service.

\subsection{Dataset and variable description}

The dataset-level record is a \code{dcat:Dataset} in HealthDCAT-AP. The project title
gives \code{dct:title}, and the identifier is built from the project identifier and an
institutional base address. The variables are described with the W3C vocabulary for CSV
on the Web (CSVW),\citep{csvw} which HealthDCAT-AP uses for this purpose. Each REDCap
form becomes a \code{csvw:Table}, and the tables of a project form one
\code{csvw:TableGroup} linked from the dataset. Each field becomes a \code{csvw:Column}
whose name, label and note come from the dictionary. The datatype follows from the field
type and its validation: an integer or number validation gives \code{integer} or
\code{decimal}, a date validation gives \code{date}, a yes--no field gives
\code{boolean}, and a coded radio or drop-down field gives \code{integer}. A checkbox
field with $n$ options is exported by REDCap as $n$ columns, and the description follows
the exported shape. Event and repeat-instance columns, which appear in the export of a
longitudinal project but not in its dictionary, are recovered from the event definitions
of the XML export. The REDCap flag \emph{Identifier?} becomes a column-level personal
data flag.

\subsection{Composition with the data management plan}

The RDA DMP Common Standard is a JSON schema with a companion ontology that specialises
the W3C DCAT vocabulary,\citep{rdadmp,cardoso2022dcso} which is also the base of
HealthDCAT-AP. Most of its dataset elements therefore have a direct counterpart in the
profile (Table~\ref{tab:dmp}). The plan is parsed and validated with \code{madmpy}.%
\citep{madmpy} A plan may describe several datasets, and the adapter selects the one
whose title matches the REDCap project.

\subsection{Terminology coding at three levels}

HealthDCAT-AP records the code systems used in a dataset with
\code{healthdcatap:hasCodingSystem}. We emit codes at three levels. At the dataset level,
\code{hasCodingSystem} lists each system that occurs in the project. At the form level,
each \code{csvw:Table} lists with \code{dct:subject} the concepts of its fields. At the
field level, each \code{csvw:Column} links with \code{skos:exactMatch} to the concept it
denotes, for instance \url{http://snomed.info/id/22298006}, with the concept's label and
the version of the code system.

REDCap has no dedicated place for terminology codes. Each field has a free-text
\emph{field annotation}, which REDCap uses for action tags such as \code{@HIDDEN}. We
carry codes there in a repeatable tag of the form
\code{@CODE=<system>|<version>|<code or expression>}, for example
\code{@CODE=SNOMED|20230731|22298006} or \code{@CODE=LOINC|2.74|4548-4}. The system token
names any code system, the version is optional, and the value is either a code or a
SNOMED CT compositional expression.\citep{snomedct} A plain code yields a concept
address; an expression is carried as text. The tag ends at the next action tag, so the
behaviour of the form is unchanged. \TBD{Decide: also read the \code{@ONTOLOGY} JSON
written by ROME, and say so here.}

The method emits codes that a project already carries and does not decide which concept
a field denotes. The tool includes a demonstration mapper that assigns SNOMED CT concepts
from field labels with a small lookup table. Its assignments are not clinically
validated, and we use it only to exercise the three levels of emission.

\subsection{Implementation}

The tool, \code{redcap2hdcat}, is written in Python and released under the Apache~2.0
licence. Each source is read by an adapter into one intermediate representation of the
project, and emitters serialise that representation as a HealthDCAT-AP record in RDF
and as an aggregated catalogue. A new source format therefore needs a new adapter and no
change elsewhere. The command-line interface has commands to build a record, to validate
it, to audit a data dictionary for design problems such as coded fields without value
sets or identifiers not flagged as such, to produce a synthetic sample distribution, and
to remove patient records from an export. Validation runs pySHACL over the published
SHACL shapes.\citep{shacl} A record built for one project is a single command:
\code{redcap2hdcat build project.xml -c institution.yaml --dmp plan.json -o record.ttl}.

\subsection{Evaluation}

\paragraph{Reference set} We took as reference the minimum metadata elements listed in
Part~B of the Annex to Implementing Regulation 2026/2098\citep{ir20262098} and mapped
each to its HealthDCAT-AP property in release \TBD{pinned release, date}.%
\citep{healthdcatap}

\paragraph{Source contribution} For one clinical research project, a longitudinal study
of rectal cancer at \TBD{institution}, we built the record in three cumulative stages:
(a) from the REDCap export alone, with no institutional settings; (b) with the
institutional settings added; and (c) with a DMP added. For each element we recorded the
stage that first supplied it. The plan used in stage (c) was written for this
evaluation from the study's documentation \TBD{confirm: the plan is the synthetic test
fixture, not the study's plan}. For the public projects described below, which have no
published plans, we repeated stages (a) and (b).

\paragraph{Conformance} We validated each record against the SHACL shapes of the pinned
release, first with the health data access body left empty, as the tool does when none
has been designated, and then with a placeholder body, to separate the violations caused
by that element from the rest.

\paragraph{Robustness} We ran the tool over \TBD{N} public REDCap projects from
\TBD{sources} (Table~\ref{tab:corpus}). They differ in input format, in the header
convention of the dictionary, and in whether they are longitudinal. The scripts and their
outputs are archived with the tool.\TBD{Zenodo DOI}

\subsection{Ethical considerations}

This work did not involve human subjects. The tool reads the structure of REDCap
projects and rejects exports that contain patient records. The clinical project was used
through a metadata-only export \TBD{confirm institutional authorisation to describe the
rectal cancer project}, and no patient data were accessed. Ethical review was therefore
not required.

\section{Results}

\subsection{Contribution of each source}

\TBD{From results/contribution.csv: number of the required elements supplied at stage
(a), (b) and (c); elements not supplied by any source and whether they need the REDCap
programming interface (number of records, number of unique individuals, temporal
coverage) or a person (description, health theme). Report the personal-data case: the
REDCap export flags \TBD{n} fields as identifiers, and the declaration of personal data
comes from the plan.} Table~\ref{tab:elements} lists the elements with the source of each.

\TBD{From results/contribution\_corpus.csv: across the public projects, the elements that
REDCap alone supplied in all of them, and those it supplied in none.}

\subsection{Conformance}

\TBD{From results/shacl\_summary.csv: violations with the access body empty and with a
placeholder; each remaining violation type and its cause.}

\subsection{Robustness}

\TBD{From results/corpus.csv: number of projects, forms and variables processed; failures
found and fixed, including the two header conventions of the data dictionary (the
download from the user interface uses \emph{Variable / Field Name}, the programming
interface uses \emph{field\_name}).} Table~\ref{tab:corpus} summarises the projects.

\section{Discussion}

\subsection{Principal findings}

\TBD{Two or three paragraphs once results are in. Expected line: the governance elements
come from outside the capture system; most of what a person supplies is set once per
institution; the plan supplies the per-study governance values; composition corrects
the personal-data declaration that REDCap alone gets wrong.}

\subsection{Relation to existing tools}

The guideline of TEHDAS2 states what each element should contain and does not address
how descriptions are produced from existing systems, and the online editor requires a
person to enter each description.\citep{tehdas2d51,hipeditor} The method described here
produces the elements that can be derived and leaves to a person only those that no
source holds. ROME and RareLink
add terminology annotations to REDCap projects as they are designed,%
\citep{meigen2026rome,graefe2025rarelink} and the method can carry such annotations into
a catalogue record. \code{sdc\_healthdcat\_ap} derives records from models of the
Semantic Data Charter and takes the governance values from a file that the data holder
writes;\citep{sdchealthdcat} the method described here reads them from a plan that the
funder usually requires already. Tools such as DMPTool and RDMO write machine-actionable
plans, and the method reads their output. Beacon~v2 lets a user find individual records
that match genomic or phenotypic criteria across federated nodes.\citep{rambla2022beacon}
An EHDS catalogue lists datasets with their access conditions, and HealthDCAT-AP is the
profile required for it.

\subsection{Implications for public health}

\TBD{Two to four sentences: data holders must describe their datasets by 2029; REDCap is
the capture system of many academic studies; deriving the description from what already
exists lowers the cost per dataset to the per-study governance values, which are already
in the plan when funders require one; catalogues with variable-level and coded
descriptions let researchers find comparable datasets across institutions before
requesting access from a health data access body.}

\subsection{Limitations}

The REDCap file formats do not carry the number of records, the number of individuals or
the temporal coverage. These need the REDCap programming interface, for which an adapter
is in development. The plan is matched to the project by title, which fails when the
titles differ; a persistent identifier shared between the plan and the project would be
more reliable. Code generation is limited to the demonstration mapper, and the quality of
automatic annotation is a separate question. Spain had not designated its health data
access body at the time of writing,\citep{hdabes} so records for Spanish data holders
cannot yet carry that element. The public projects come from consortia that publish
their instruments for reuse, so they are not a sample of typical projects; we use them
to test robustness and not to measure dictionary quality. A study of dictionary quality
over the projects of an institutional REDCap instance is in preparation.

\section{Conclusions}

\TBD{Qualitative, no digits: a HealthDCAT-AP description can be composed from the
capture system, the plan and institutional settings; which elements each supplies; what
remains for the programming interface and for a person.}

\section*{Software availability}
Source code: \TBD{repository URL} (Apache~2.0). Archived release: \TBD{Zenodo DOI}.
Evaluation scripts and results: \TBD{DOI}.

\section*{Acknowledgments}
\TBD{Technical assistance, if any.}

\section*{Funding}
This work was carried out within the project IMPaCT-Data~2 (PMP24/00024,
ISCIII-AES-2024/003581), funded by the Instituto de Salud Carlos~III (ISCIII) under the
\emph{Misiones 2024} call for Personalised Precision Medicine, and co-funded by the
European Union through NextGenerationEU and the Recovery and Resilience Facility.
\TBD{Confirm the funder's required wording.}

\section*{Conflict of Interest}
None declared. \TBD{Confirm for every author; ICMJE forms.}

\newpage
\bibliographystyle{unsrtnat}
\bibliography{refs}

\newpage
\section*{Figure captions}

\noindent\textbf{Figure 1.} Architecture of the method. Adapters read the REDCap
dictionary or project export, the machine-actionable DMP and the institutional settings
into one intermediate representation; the coding layer attaches terminology concepts;
emitters write the HealthDCAT-AP record with its CSVW variable description.
\TBD{Draw as black-and-white line art; submit as a separate EPS/TIFF file.}
\label{fig:architecture}

\newpage
\section*{Tables}

\begin{table}[h]
\caption{Minimum metadata elements of Implementing Regulation 2026/2098 and the source
that supplies each in the rectal cancer project.}
\label{tab:elements}
\small
\TBD{Generated from eval/results/contribution.csv: element; HealthDCAT-AP property;
condition in the Annex; source (REDCap / institutional settings / DMP / REDCap programming
interface / person).}
\end{table}

\begin{table}[h]
\caption{Elements of the RDA DMP Common Standard read into HealthDCAT-AP.}
\label{tab:dmp}
\small
\begin{tabular}{@{}ll@{}}
\toprule
DMP element & HealthDCAT-AP property \\
\midrule
\code{dataset.description} & \code{dct:description} \\
\code{contributor} (role, ORCID) & \code{dct:creator} \\
\code{contact} & \code{dcat:contactPoint} \\
\code{distribution.license} & \code{dct:license} \\
\code{project.funding.grant\_id} & \code{dct:isReferencedBy} \\
\code{dataset.personal\_data} & \code{dpv:hasPersonalData} \\
\code{dataset.keyword} & \code{dcat:keyword} \\
\code{dataset.language} & \code{dct:language} \\
\code{dataset.issued}, \code{modified} & \code{dct:issued}, \code{dct:modified} \\
\bottomrule
\end{tabular}
%% VERIFY this table against src/redcap2hdcat/adapters/dmp.py before submission.
\end{table}

\begin{table}[h]
\caption{Public REDCap projects used to test robustness.}
\label{tab:corpus}
\small
\TBD{Generated from eval/results/corpus.csv: source, domain, input format, header
convention, forms, fields, longitudinal.}
\end{table}

\end{document}
